\section*{Chapitre \ref{ch:s1:flows} --- Flux}

\begin{solution}{exo:s1:flows:1}
$0.2 \times 0.05 = 1.0~\%$ des actions du titre à acheter~; avec un impact de 1.5, une poussée de 1.5~\% sur le trimestre.
\end{solution}

\begin{solution}{exo:s1:flows:2}
La moitié (50~\%) après 126 jours, un quart après 252~: les trois quarts de la poussée se sont inversés en un an.
\end{solution}

\begin{solution}{exo:s1:flows:3}
Déposées le dernier jour permis, les positions du 31 mars ont 45 jours à leur réception (mi-mai)~; elles restent les plus récentes jusqu'à l'arrivée des positions de juin
à la mi-août, date à laquelle elles ont environ quatre mois et demi.
\end{solution}

\begin{solution}{exo:s1:flows:4}
Les fonds qui ont bien performé détiennent des actions performantes~; les actions synthétiques ont des dérives persistantes, si bien que les actions détenues par les bons fonds continuent de monter. Les flux poursuivent les
rendements des fonds et sont négociés sur ces mêmes actions, de sorte que les \omterm{def:s1:flows:fit}{transactions induites par les flux} et les rendements passés s'alignent même sans pression (une pente en coupe transversale de 0.31).
\end{solution}

\begin{solution}{exo:s1:flows:5}
Un fonds aux bons rendements reçoit des entrées, achète davantage de ce qu'il détient et fait monter ces cours~; ce qui relève ses rendements suivants, qui attirent de nouvelles entrées.
Sa performance persiste à cause des flux qu'elle attire, et non de son talent~: c'est l'explication de Lou.
\end{solution}

\begin{solution}{exo:s1:flows:6}
La transaction ne peut commencer qu'une fois connues les opérations du trimestre, 45 jours après la fin de celui-ci, moment où une partie de l'inversion a déjà eu lieu~; ce qui reste est
faible et se heurte aux dérives que poursuivaient les flux, qui maintiennent à la baisse les actions vendues. L'\omterm{def:s1:earnings:event}{étude d'événement} mesure l'inversion à partir du
début du trimestre~; le livre n'intervient que bien plus tard.
\end{solution}

\begin{solution}{exo:s1:flows:7}
Le décile supérieur moins le décile inférieur gagne 4.9~\% sur le trimestre et 69.0~\% sur trois ans. Les fonds orientés vers un style rémunéré en perçoivent la prime, attirent des entrées et
achètent ce style, de sorte que le signal de flux classe les actions selon leur style. Pour tester un véritable signal de flux, il faut régresser ses rendements sur les facteurs de style et le momentum, ou
le construire à partir de flux résiduels sur les rendements de style du fonds.
\end{solution}

\begin{solution}{exo:s1:flows:8}
Les positions de fin de trimestre sont déposées jusqu'à 45 jours plus tard, de sorte qu'une négociation à la date de fin de trimestre utilise une information qui n'est pas encore publique~; les flux sont également publiés avec un retard.
Le backtest honnête négocie aux dates de dépôt et contrôle le momentum et les styles.
\end{solution}

\begin{solution}{pb:s1:flows:1}
\begin{enumerate}
\item Transactions imposées par les flux des fonds (part de détention multipliée par le flux, sommée sur les fonds)~; la partie d'un mouvement de cours causée par une demande non informée, censée s'inverser.
\item Les fonds à fortes sorties vendent leurs positions existantes, créant une pression sur les actions communément détenues~; négocier contre eux rapporte des rendements significatifs~; les
transactions forcées sont prévisibles.
\item Un impact sur les prix significatif mais temporaire des \omterm{def:s1:flows:fit}{transactions induites par les flux}~; sa partie attendue prévoit les rendements de l'année suivante, puis s'inverse~; il peut expliquer la
persistance de la performance des fonds, l'effet smart money et une partie du momentum.
\item Les institutions sont contraintes, la demande du marché est inélastique, et 1~\$ investi élève sa valeur d'environ 5~\$.
\item Deux cents fonds de soixante actions tirées au hasard, un quart remplacé à chaque trimestre, détenant 20~\% de chaque action~; des flux de 5~\% par écart type du rendement
passé glissant plus un bruit de 3~\%~; une poussée de 1.5 par unité de transaction induite par les flux sur le trimestre, s'inversant avec une demi-vie de 126 jours.
\item La poussée relève les rendements des fonds qui détiennent les actions poussées, lesquels attirent alors davantage de flux.
\item 1.05 sur le trimestre et $-0.69$ l'année suivante avec pression~; 0.31 et 0.32 sans.
\item 66~\% (75~\% de la poussée de fin de trimestre due au seul déclin programmé).
\item Flux prévus à partir des rendements passés glissants des fonds multipliés par les positions déposées~; ratio de Sharpe de 0.44 après coûts.
\item Elle intervient 45 jours après le trimestre, trop tard pour une bonne part de l'inversion, et se heurte aux dérives~; $-0.35$.
\item Des transactions de paniers mécaniques par arbitrage~: plus de volatilité et autocorrélation négative dans les actions qu'elles détiennent.
\item Rapports trimestriels de positions des gérants supérieures à 100 millions de dollars en titres 13(f), déposés dans un délai de 45 jours.
\item \textbf{Résultat nommé.} Chaque unité de transaction induite par les flux déplace le rendement logarithmique du trimestre de 1.05 (dont 0.31 relève de la dérive programmée que poursuivent les flux), et 66~\%
du mouvement du trimestre s'inverse dans l'année qui suit.
\item Le décile supérieur moins le décile inférieur gagne 1.74~\% sur le trimestre, culmine à 1.91~\% juste après, revient à zéro après environ un an et termine à $-2.15~\%$ après
trois ans.
\item Avec des fonds orientés vers un style, le signal devient un pari sur le style~: 69.0~\% sur trois ans, sans jamais s'inverser.
\item Actifs et flux mensuels des fonds, créations quotidiennes d'ETF, et rapports au niveau des fonds sur les changements de portefeuille.
\item Il faut régresser ses rendements sur le momentum et les facteurs de style, ou résidualiser les flux sur les rendements de style des fonds.
\item Le livre de flux attendus, avec le modèle de flux réestimé à mesure que les flux évoluent.
\item Une connaissance plus rapide des ventes forcées~: données de flux quotidiennes ou nouvelles des liquidations.
\item La liquidité, pour les investisseurs qui doivent négocier.
\end{enumerate}
\end{solution}

\begin{solution}{iq:s1:flows:1}
Parce que les fonds doivent négocier leurs flux dans les actions qu'ils détiennent, et que la contrepartie de cette demande n'est pas parfaitement élastique~; le cours bouge jusqu'à ce que quelqu'un soit
suffisamment rémunéré pour la prendre, puis revient à mesure que la demande est absorbée.
\end{solution}

\begin{solution}{iq:s1:flows:2}
Les positions issues des dépôts 13F et des dépôts de portefeuille des fonds, les flux issus des variations d'actifs des fonds nets des rendements, chacun daté par sa publication~; les \omterm{def:s1:flows:fit}{transactions induites par les flux} sont la somme sur
les fonds de la part de détention multipliée par le flux. Le délai va de quelques semaines à plusieurs mois~: positions jusqu'à 45 jours après le trimestre pour le 13F, flux mensuels ou trimestriels.
\end{solution}

\begin{solution}{iq:s1:flows:3}
Régresser ses rendements long--short sur le momentum et les facteurs de style, trier au sein des déciles de momentum, et construire la mesure de flux à partir de flux résiduels sur le style
du fonds et ses rendements passés~; un signal qui résiste relève bien des flux.
\end{solution}

\begin{solution}{iq:s1:flows:4}
Estimer ce qu'il doit vendre et à quelle vitesse, ne négocier contre lui qu'une fois que ses ventes ont poussé les cours sous la valeur, et tenir jusqu'à l'inversion. Se méfier des autres
front-runners, du fonds qui vend les mêmes titres que d'autres sous tension, et de la question de savoir si la liquidation reflète une information.
\end{solution}

\begin{solution}{iq:s1:flows:5}
Ventes corrélées~: lorsque ces fonds font face à des sorties, l'action est vendue avec ses autres détentions communes, indépendamment de ses propres nouvelles, de sorte que son risque inclut le
risque de flux des fonds.
\end{solution}

\begin{solution}{iq:s1:flows:6}
La poussée restante $s$ jours après le trimestre vaut $\lambda F\,2^{-s/h}$. Acheter (pour $F < 0$) en $L$ et vendre en $L + H$ rapporte
$-\lambda F\,2^{-L/h}(1 - 2^{-H/h})$, positif pour les actions vendues, avant coûts $2c$~; la transaction est rentable lorsque $|\lambda F|\,2^{-L/h}(1 - 2^{-H/h}) > 2c$. Un long délai
$L$ la réduit exponentiellement, ce qui explique pourquoi le livre des ventes forcées a besoin de données rapides.
\end{solution}
